\section{Security Context and Threat Model}

\paragraph{Capability mediation.}
Capability systems replace ambient authority with explicit, attenuable references. KeyKOS and EROS made this discipline an operating-system primitive; Capsicum brought capability mode to UNIX; seL4 connected a microkernel authority model to machine-checked functional correctness \cite{hardy1985keykos,shapiro1999eros,watson2010capsicum,klein2009sel4}. Object-capability work explains authority as reachability and attenuation rather than membership in a global access-control list \cite{miller2006composition,murray2010ocap}. Chio applies that discipline to agent tool calls: the untrusted agent supplies a signed capability, while a trusted kernel validates it and mediates the tool server.

\paragraph{Agent isolation and cross-domain authority.}
IsolateGPT confines agent applications behind information-flow gates; SAGA distributes user-controlled authorization tokens; recent agent-security work identifies confused-deputy and authority-boundary failures as a central systems problem \cite{isolategptNDSS2025,sagaNDSS2025,agenticFoundations2025}. These systems primarily protect a local principal. Bilateral admission begins when the evidence crosses an administrative boundary and the receiver must decide whether a foreign action belongs to its own authorized history.

\paragraph{Signed statements are necessary but incomplete.}
DSSE authenticates a payload and predicate type, while in-toto and SLSA bind subjects to supply-chain provenance \cite{dsse,torres2019intoto,slsa}. A valid provenance statement does not establish receiver authorization for a tool call. The receiver additionally needs freshness, request and outcome binding, participant and treaty binding, continuation state, and its own policy verdict. Chio retains DSSE's canonical signed envelope but defines an action-admission predicate rather than a build-provenance predicate.

\paragraph{Trusted components.}
We trust the receiving kernel, its cryptographic and canonical-JSON libraries, and its locally provisioned stores. The stores contain expected peer keys, treaty scope, ladder manifests, leases, governance receipts, revocation state, continuations, lineage, and resolved receipts. We assume Ed25519 and SHA-256 behave as specified and that the receiver classifies federated origin correctly. A compromised receiver can authorize anything and lies outside the property proved here.

\paragraph{Adversary.}
The adversary controls the agent request and may present arbitrary metadata and syntactically valid signed objects. It may replay old material; mix artifacts from different invocations; alter the treaty, request, outcome, receipt, lineage, continuation, or predicate type; omit a signature or governance record; use stale leases; duplicate a signer key identifier; or submit a correctly signed unanimous denial. It may control a remote vendor kernel and any request-carried trust data. The receiver must deny these cases before its tool runs.

The adversary may also control two private keys that the receiver has separately authorized. Cryptography cannot distinguish this condition from two administrative principals. We therefore test rejection of an identical repeated key but record distinct-controller independence as assumption \textsc{PS-A-01}. TEE-backed key custody or external organizational governance could strengthen that assumption, but neither is part of the evaluated construction.

\paragraph{Security objective.}
Let $q$ be a federated tool request and $A_R(q)$ the receiving kernel's admission decision. The implementation objective is:
\[
\begin{split}
A_R(q)=\mathsf{allow}\ \Longrightarrow\ &
  \mathsf{coreAllows}(q)\ \land\\
& \mathsf{receiverEvidenceAccepts}(q).
\end{split}
\]
The second predicate means that receiver-resolved treaty evidence is fresh,
mutually bound, strictly verified, and accepted by receiver policy.
If any premise fails, the receiver returns a denial before invocation. This is a local safety property. It neither forces the sender to record the same history nor establishes the truth of claims made by a compromised but authorized signer.
